% généré par scripts/04_generated.py
\begin{longtable}{@{}r r l p{0.6\linewidth}@{}}\caption{Datations contestées, réserves et divergences attribuées à l’EDB. L’année de classement est conservée dans tous les cas.}\label{tab:contestees}\\\toprule N° & Année & Statut & Motif\\\midrule\endfirsthead\toprule N° & Année & Statut & Motif\\\midrule\endhead\bottomrule\endfoot
40 & 353 & contestée & {\footnotesize Magnence III et Décence II (353) selon la lecture retenue ; l’EDB, citant CLRE p. 636, date de 352 en supposant une erreur du lapicide.} \\
52 & 358 & contestée & {\footnotesize Da[tiano] (358) selon Silvagni ; la CLRE (p. 631, citée par l’EDB) rappelle que Dalmatius, consul en 333, commence aussi par DA. Si 333, la notice sort du corpus.} \\
141 & 370 & contestée & {\footnotesize Valentinien III : Valentinien Ier (370) ou Valentinien II (387) ; la CLRE (p. 644, citée par l’EDB) note « Year doubtful: 387 possible ».} \\
165 & 372 & contestée & {\footnotesize L’année 372 repose sur la seule séquence « dest », interprétée par Ferrua comme [Mo]dest[o] ; l’EDB date « 300-399 ».} \\
206 & 370 & EDB fautive & {\footnotesize L’EDB date de 377 ; la formule Valentiniano III et Valente Augg. ne convient qu’à 370 : erreur probable de saisie de l’EDB.} \\
220 & 381 & contestée & {\footnotesize Seul …erio est lu ; l’EDB signale d’autres restitutions possibles (Val]erio, Gal]erio, Neot]erio, Ast]erio) : 381 ou 390.} \\
278 & 383 & EDB fautive & {\footnotesize L’EDB date de 385 ; Merobaude II et Saturnino = 383, et la fiche voisine EDB8750 (même consulat) est datée de 383 : incohérence interne probable de l’EDB.} \\
311 & 388 & contestée & {\footnotesize Maximo cons(ule) : consulat de Magnus Maximus (388) ; l’EDB date de 389 et signale une transcription différente de l’ICVR.} \\
314 & 390 & contestée & {\footnotesize Second texte du monument ICUR IV 11117 : seul d(omino) n(ostro) V[…] est lu ; 390 est une restitution partagée par l’EDCS et l’EDB, que la note de l’EDB juge contradictoire avec l’ordre des décès.} \\
328 & 381 & contestée & {\footnotesize Seul …terio est lu : Eucherio (381) ou Neoterio (390).} \\
395 & 396 & réserve & {\footnotesize Itération d’Honorius (IIII transmis) à confronter au collège de 396 (Arcadius IV et Honorius III).} \\
411 & 397 & réserve & {\footnotesize Collège Caesarius et Maximus ; les fastes donnent pour 397 Caesarius et Atticus.} \\
428 & 398 & réserve & {\footnotesize Post-consulat d’Honorius IV et Eutychianus (399 ?) ou consulat (398) ; le jour de la semaine (samedi) convient au 18 décembre 398.} \\
431 & 399 & contestée & {\footnotesize La CLRE (p. 661, citée par l’EDB) attribue le texte à 398 ; 399 suppose un post restitué que l’EDB juge trop long.} \\
434 & 397 & contestée & {\footnotesize Post-consulat d’Arcadius IV : 397 ; l’EDB date de 399.} \\
444 & 400 & réserve & {\footnotesize Itération [iter(um)] de Stilichon restituée ; compatibilité avec 400 à contrôler.} \\
466 & 402 & réserve & {\footnotesize Paire Arcadius–Théodose à contrôler sur les fastes de 402.} \\
\end{longtable}
